\scanpage{99}\mspage{90}
\deleted{Si \unsure{donc} on}

\uline{Corollaire} Soit $\unsure{\mathscr{T}|\mathcal{T}}\unsure{\subset|\subseteq}\mathrm{Drap}^{\unsure{*|\times}}(L)$ un ens \unsure{de|des} drapeaux non vides. Pour que les figures $F_T$ ($T\in\mathscr{T}$) soient deux à deux disjointes, il f. et s. que $\mathscr{T}$ soit un drapeau de l'ens ordonné $\mathrm{Drap}^{*}(L)$.

Dans ce cas, on a donc une figure somme
\[
 F_{\mathscr{T}}\overset{\text{def}}{=}\bigcup_{T\in\mathscr{T}}F_T
\]
% slanted
\margin{Les $F_T$}
on retrouve $F_T$ comme cas particulier
\[
 F_T=F_{\{T\}}
\]

La figure vide correspond au cas $\mathscr{T}=\varnothing$.

\uline{Théorème} L'application
\[
\begin{aligned}
 &\mathscr{T}\longrightarrow F_{\mathscr{T}}\\
 &\mathrm{Drap}(\mathrm{Drap}^{*}(L))\longrightarrow
 \mathrm{Fig}\,\mathrm{fin}(\unsure{M|\mathcal M})=\unsure{\mathcal{F}|\mathscr{F}}
\end{aligned}
\]
est bijective.

Indication de la démonstration. Pour une figure donnée, on récupère les $F_T$ comme les ``composantes connexes'' (\deleted{\unsure{au|en}} \inserted{dans un} sens combinatoire \unsure{'}évident) de\deleted{\unsure{s}} $F$. Il \deleted{\unsure{f|y}} reste à voir que les composantes connexes d'une figure sont bien de la forme $\deleted{\unsure{E|F}}F_T$, puis utiliser le lemme précédent.
% archive sheet number: 99
